\documentclass[11pt]{article}

\usepackage[margin=1in]{geometry}
\usepackage{times}
\usepackage{microtype}
\usepackage{booktabs}
\usepackage{hyperref}
\usepackage{xcolor}
\usepackage{enumitem}
\usepackage{listings}

\hypersetup{colorlinks=true, linkcolor=black, citecolor=black, urlcolor=blue}

\lstset{
  basicstyle=\ttfamily\small,
  breaklines=true,
  columns=fullflexible,
  frame=single,
  framesep=4pt,
  xleftmargin=4pt,
  keepspaces=true,
}

\title{\textbf{Ledgerkeep}\\[2pt]
\large A Guardrailed Autonomous Operations Agent as a Decentralized AI Service}
\author{Dipankar Sarkar \\ \small doom2quake}
\date{August 2026 \\ \small Version 0.1.0 (Milestone 1)}

\begin{document}
\maketitle

\begin{abstract}
Anomaly detection is a solved commodity; the unsolved step is the one after the
alert. An operations agent that a team would let near its systems must be
trusted to investigate, attribute a moved number to a specific change, quantify
the loss, and propose the exact fix, \emph{without} being trusted to execute
anything a human has not approved. That gap is almost entirely guardrails,
durable memory, correct routing, and an honest report of what the agent actually
did. Ledgerkeep is that layer, packaged as a decentralized AI service on the ASI
/ Fetch.ai stack: a stable, typed request and response, and a verifiable audit
object returned with every call. This paper documents the milestone-1
deliverable: the service contract, the guardrailed investigation loop, and the
adapter seam that makes the funder's technology load-bearing while keeping the
whole system runnable offline and keyless. We report reproducible test counts and
state plainly what has not been built.
\end{abstract}

\section{The problem}

Consider the incident Ledgerkeep is named for. A payments configuration is
pushed at 14:03. It routes one region's card traffic to a gateway that declines
most authorisations. Order volume still looks normal, because customers keep
trying to buy; completed, settled revenue quietly falls, and nobody notices
until the next morning. The person who suffers this is the on-call operator,
paged at 03:00, who must reconstruct from a dashboard that shows only the symptom
which of forty deploys that day moved a number, how much money has already been
lost, and what one line to change. That reconstruction is slow, it happens under
pressure, and it is where the cost lives.

The reason this stays unsolved is not detection. Trailing-window outlier
detection is textbook~\cite{grubbs_1969_outliers}. The unsolved part is the step
\emph{after} the alert: an agent trusted to investigate and propose, but not to
act. The market is full of demonstration agents that will happily run generated
SQL against production or act on model text that contains an
injection~\cite{owasp_llm_top10_2025}. The distance between a demo agent and one
an operator would let near their systems is guardrails, durable memory, correct
routing, and an honest report. That is the layer nobody wants to build twice.

\section{Why a decentralized service, and why ASI}

ASI Deep Funding funds decentralized AI \emph{services}: agents that others can
discover, call, and pay for, while the builder keeps ownership~\cite{asi_deepfunding_2026}.
An operations agent that a team must trust is exactly the kind of component that
benefits from being a standalone, independently callable service with a
published, auditable contract, rather than a black box buried inside one
company's stack.

Ledgerkeep maps onto that thesis directly. The agent is a marketplace service: a
stable, typed request and response, a guarantee about what it will and will not
do, and an audit object returned with every call. On the ASI / Fetch.ai stack a
service is a \texttt{uAgents} agent registered on
Agentverse~\cite{fetchai_uagents_2026}; the SingularityNET marketplace is where
such a service is discovered and monetized~\cite{singularitynet_marketplace_2026}.
The Model Context Protocol~\cite{mcp_spec_2025} is the interoperability seam: the
same agent that answers a marketplace call can be consumed as a tool by another
agent, which is the composition story the Alliance is built around.

\section{The service contract}

Milestone 1 delivers Ledgerkeep as a callable service, so the interface is the
product. The contract is transport-agnostic: plain typed objects with an explicit
JSON round-trip, mapped onto a uAgents message by the marketplace adapter, onto
an MCP tool by the MCP bridge, or exchanged directly by a local caller. One
contract, many transports.

\paragraph{Request.} \texttt{InvestigationRequest} is typed and all-defaults, so
the simplest possible call is \texttt{InvestigationRequest()}. Fields are the KPI
to watch, the z-score threshold, a free-text note recorded for the audit trail
but never executed, and a caller correlation id echoed back on the response.

\paragraph{Response.} \texttt{InvestigationResponse} carries the anomaly, the
attribution, the quantified impact, the proposed fix, a verification note, and a
\texttt{delivery\_mode} that is always present and reads \texttt{offline\_fixture}
when the figures came from the in-repo ledger rather than a live source, so a stub
is never mistaken for a live figure. Two invariants are enforced by the schema
itself: the proposed fix is always returned with \texttt{approved = false} (the
merge stays with a human), and the response cannot be constructed without an
audit object.

\paragraph{Audit object.} \texttt{AuditObject} is returned on \emph{every} call.
It records the run id, the ordered guardrail decisions by name, the classified
domain, the recurrence record, and a SHA-256 \texttt{content\_hash} over the
canonical JSON of the decision chain. A caller can recompute the hash from the
visible fields and confirm the audit trail was not edited. This is the surface
that milestone 3 anchors to a public testnet; here it is computed and returned,
so the tamper-evidence primitive already exists.

\begin{lstlisting}
audit.run_id        = run-20260825T204713-a75959
audit.status        = acted
audit.domain        = finance
audit.guardrails    = [CONTENT_SAFETY:clean, DOMAIN_ROUTER:classified,
                       ACTION_LIMITER:allowed]
audit.content_hash  = 7d735e4ee91547f25deb2bd262182fd85439c0a87ed79c...
\end{lstlisting}

\section{The guardrailed loop}

\texttt{run\_service} takes a request and runs a deterministic, offline loop
against a fixture ledger, enforcing every guardrail from the vendored control
plane and recording each decision by name on a durable run document.

\begin{enumerate}[leftmargin=1.4em, itemsep=1pt]
  \item \textbf{Detect} the settled-revenue drop by trailing z-score.
  \item \textbf{Attribute} it to a specific change by a deterministic rule over
        the change-event log. The attribution text is screened by
        \texttt{CONTENT\_SAFETY} before anything downstream trusts it; a flagged
        run is quarantined with no proposed fix, so a poisoned diagnosis cannot
        talk the agent into a remediation.
  \item \textbf{Route} the incident through a keyword domain classifier so the
        audit trail records the classified domain.
  \item \textbf{Quantify} the loss to the dollar.
  \item \textbf{Propose} the one-line fix, returned unapproved. Recording the
        proposal passes the \texttt{ACTION\_LIMITER} rate cap.
  \item \textbf{Verify} by re-reading the metric after a fixture failover, as the
        next live cycle would after a real remediation.
\end{enumerate}

The bound on autonomy is three named guardrails, each tested without a model. A
security incident that were also a finance incident is routed to security first,
because misrouting it would lose the incident~\cite{nygard_2018_release}. Nothing
in this loop reaches the network.

\section{The adapter seam}

The funder's technology is load-bearing behind a seam. When a seed phrase is set
and the \texttt{uagents} SDK is installed, on a test network, the marketplace
adapter builds a real \texttt{uagents.Agent}, registers a handler that decodes an
incoming message into an \texttt{InvestigationRequest}, runs the identical
\texttt{run\_service}, and replies with the response. Otherwise an in-process
transport answers the exact same contract, so the service is callable and
testable with no SDK and no credentials. The live handler is a one-line funnel
into the same call the offline path uses, so the two transports cannot diverge in
behaviour. All on-chain interaction is testnet only for the duration of the
grant; a non-test network is refused.

A \texttt{describe()} manifest is produced from the same source of truth the
adapter uses, reporting the service title, network, contract version, the request
and response model names, and which transport is active, so a published
description can never drift from the running service.

\section{Evaluation}

The suite is hermetic: it pins the environment offline and in-memory and cuts the
socket, HTTP, and subprocess layers before the package is imported, so nothing in
the suite can reach an external service even by accident. Tests that would
exercise an integration inject their own fake transport.

\begin{center}
\begin{tabular}{lr}
\toprule
\textbf{Surface} & \textbf{Tests} \\
\midrule
Service contract (round-trip, audit hash, invariants) & 10 \\
Guardrailed loop (detect / attribute / quantify / propose / verify) & 7 \\
Marketplace adapter (transports, manifest, testnet refusal) & 10 \\
Fixture ledger (deterministic scenario, failover) & 7 \\
Vendored \texttt{agent\_core} control plane & 48 \\
\midrule
\textbf{Total (all passing, offline)} & \textbf{82} \\
\bottomrule
\end{tabular}
\end{center}

Every assertion about the loop is read off the returned audit object: a fact
about what the agent did, not a claim about what it said. Two tests confirm the
guardrails bite: a poisoned attribution quarantines the run with no fix, and a
zero-cycle action cap blocks the proposal and records the block.

\section{Honest limits}

Milestone 1 is a callable, tested, self-contained service. It is not a
deployment. \textbf{There are no users, no revenue, no partnerships, and no
external audit.} There is \textbf{no live ASI marketplace publication} in this
repo: the live \texttt{uAgents} path is present behind the seam but has never been
exercised against Agentverse in this environment. \textbf{There is no mainnet},
and there will be none during the grant; the audit-hash anchoring is milestone 3
and is off by default. The content-safety and read-only-SQL screens in the
control plane are deny-list text screens, not parsers; they are defence in depth
and the real boundary for any generated query is a read-only credential with a
server-side cost cap. These limits are stated in \texttt{docs/LIMITATIONS.md} so
diligence finds nothing that was not disclosed.

\section{Conclusion}

The durable output is not one agent, it is the safety layer under it. The
vendored \texttt{agent\_core} control plane is MIT-licensed and separable, so
other ASI builders can vendor the exact tested guardrails rather than
re-implementing them. The service contract and the audit object are a worked
reference for how to publish and compose a guardrailed agent service on the
stack. Milestone 1 ships that reference, running and tested, keyless.

\bibliographystyle{plain}
\bibliography{references}

\end{document}
